% GENERATED FILE. Edit canonical lessons, then run npm run generate:books.
% canonical-source-hash: fnv1a64:eb8b74f9913f135a
% canonical-lessons: ML-C266-totuka, ML-C266-katikkuka, ML-C266-cavaykkuka, ML-C266-vilunnuka, ML-C266-manakkuka

\chapter{To touch, To bite, To chew, To swallow, To smell}
\label{ch:ml-to-touch-to-bite-to-chew-to-swallow-to-smell}
\hlchaptermodality{266}

\begin{chapteropening}
\textbf{By the end of this chapter:} \emph{I can say to touch, to bite, to chew, to swallow and to smell.}

Complete the last lesson of chapter 266: I can say to touch, to bite, to chew, to swallow and to smell.
\end{chapteropening}

\section[\texorpdfstring{\ml{തൊടുക} (toṭuka)}{തൊടുക (toṭuka)}]{\ml{തൊടുക} (toṭuka) — to touch}

\label{lesson:ML-C266-totuka}



\begin{quote}
Before the new one: say the Malayalam for to tie, then the Malayalam for to measure.
\end{quote}

\subsection*{You'll want to know: \ml{തൊടുക}}
\textbf{\ml{തൊടുക}} — \emph{toṭuka} — \textquotedblleft{}to touch\textquotedblright{}.

\begin{grammarlens}[title={What you've built}]
One more word for this chapter.
\end{grammarlens}

\subsection*{Your turn}
\begin{itemize}
\raggedright
  \item \emph{Say it:} \emph{toṭuka}
  \item \emph{Say it:} \emph{toṭuka}, once more
  \item \emph{Say it:} say \emph{toṭuka}
  \item \emph{Recall:} say the Malayalam for to tie, then the Malayalam for to measure, then say \emph{toṭuka} again
\end{itemize}

\begin{tcolorbox}[breakable,colback=teal!4,colframe=teal!35!black,title={Before you move on}]
What does \ml{തൊടുക} mean? (\textquotedblleft{}To touch\textquotedblright{}.) Say it once more.
\end{tcolorbox}
\section[\texorpdfstring{\ml{കടിക്കുക} (kaṭikkuka)}{കടിക്കുക (kaṭikkuka)}]{\ml{കടിക്കുക} (kaṭikkuka) — to bite}

\label{lesson:ML-C266-katikkuka}



\begin{quote}
Before the new one: say the Malayalam for to measure, then the Malayalam for to touch.
\end{quote}

\subsection*{You'll want to know: \ml{കടിക്കുക}}
\textbf{\ml{കടിക്കുക}} — \emph{kaṭikkuka} — \textquotedblleft{}to bite\textquotedblright{}.

\begin{grammarlens}[title={What you've built}]
One more word for this chapter.
\end{grammarlens}

\subsection*{Your turn}
\begin{itemize}
\raggedright
  \item \emph{Say it:} \emph{kaṭikkuka}
  \item \emph{Say it:} \emph{kaṭikkuka}, once more
  \item \emph{Say it:} say \emph{kaṭikkuka}
  \item \emph{Recall:} say the Malayalam for to measure, then the Malayalam for to touch, then say \emph{kaṭikkuka} again
\end{itemize}

\begin{tcolorbox}[breakable,colback=teal!4,colframe=teal!35!black,title={Before you move on}]
What does \ml{കടിക്കുക} mean? (\textquotedblleft{}To bite\textquotedblright{}.) Say it once more.
\end{tcolorbox}
\section[\texorpdfstring{\ml{ചവയ്ക്കുക} (cavaykkuka)}{ചവയ്ക്കുക (cavaykkuka)}]{\ml{ചവയ്ക്കുക} (cavaykkuka) — to chew}

\label{lesson:ML-C266-cavaykkuka}



\begin{quote}
Before the new one: say the Malayalam for to touch, then the Malayalam for to bite.
\end{quote}

\subsection*{You'll want to know: \ml{ചവയ്ക്കുക}}
\textbf{\ml{ചവയ്ക്കുക}} — \emph{cavaykkuka} — \textquotedblleft{}to chew\textquotedblright{}.

\begin{grammarlens}[title={What you've built}]
One more word for this chapter.
\end{grammarlens}

\subsection*{Your turn}
\begin{itemize}
\raggedright
  \item \emph{Say it:} \emph{cavaykkuka}
  \item \emph{Say it:} \emph{cavaykkuka}, once more
  \item \emph{Say it:} say \emph{cavaykkuka}
  \item \emph{Recall:} say the Malayalam for to touch, then the Malayalam for to bite, then say \emph{cavaykkuka} again
\end{itemize}

\begin{tcolorbox}[breakable,colback=teal!4,colframe=teal!35!black,title={Before you move on}]
What does \ml{ചവയ്ക്കുക} mean? (\textquotedblleft{}To chew\textquotedblright{}.) Say it once more.
\end{tcolorbox}
\section[\texorpdfstring{\ml{വിഴുങ്ങുക} (viḻuṅṅuka)}{വിഴുങ്ങുക (viḻuṅṅuka)}]{\ml{വിഴുങ്ങുക} (viḻuṅṅuka) — to swallow}

\label{lesson:ML-C266-vilunnuka}



\begin{quote}
Before the new one: say the Malayalam for to bite, then the Malayalam for to chew.
\end{quote}

\subsection*{You'll want to know: \ml{വിഴുങ്ങുക}}
\textbf{\ml{വിഴുങ്ങുക}} — \emph{viḻuṅṅuka} — \textquotedblleft{}to swallow\textquotedblright{}.

\begin{grammarlens}[title={What you've built}]
One more word for this chapter.
\end{grammarlens}

\subsection*{Your turn}
\begin{itemize}
\raggedright
  \item \emph{Say it:} \emph{viḻuṅṅuka}
  \item \emph{Say it:} \emph{viḻuṅṅuka}, once more
  \item \emph{Say it:} say \emph{viḻuṅṅuka}
  \item \emph{Recall:} say the Malayalam for to bite, then the Malayalam for to chew, then say \emph{viḻuṅṅuka} again
\end{itemize}

\begin{tcolorbox}[breakable,colback=teal!4,colframe=teal!35!black,title={Before you move on}]
What does \ml{വിഴുങ്ങുക} mean? (\textquotedblleft{}To swallow\textquotedblright{}.) Say it once more.
\end{tcolorbox}
\section[\texorpdfstring{\ml{മണക്കുക} (maṇakkuka)}{മണക്കുക (maṇakkuka)}]{\ml{മണക്കുക} (maṇakkuka) — to smell}

\label{lesson:ML-C266-manakkuka}



\begin{quote}
Before the new one: say the Malayalam for to chew, then the Malayalam for to swallow.
\end{quote}

\subsection*{You'll want to know: \ml{മണക്കുക}}
\textbf{\ml{മണക്കുക}} — \emph{maṇakkuka} — \textquotedblleft{}to smell\textquotedblright{}.

\begin{grammarlens}[title={What you've built}]
One more word for this chapter.
\end{grammarlens}

\subsection*{Your turn}
\begin{itemize}
\raggedright
  \item \emph{Say it:} \emph{maṇakkuka}
  \item \emph{Say it:} \emph{maṇakkuka}, once more
  \item \emph{Say it:} say \emph{maṇakkuka}
  \item \emph{Recall:} say the Malayalam for to chew, then the Malayalam for to swallow, then say \emph{maṇakkuka} again
\end{itemize}

\begin{tcolorbox}[breakable,colback=teal!4,colframe=teal!35!black,title={Before you move on}]
What does \ml{മണക്കുക} mean? (\textquotedblleft{}To smell\textquotedblright{}.) Say it once more.
\end{tcolorbox}
